\chapter{本教程使用的单片机开发板与BSP框架的构建}

本附录介绍教程各章使用的开发板，以及“如何用 STM32 HAL 从零搭建一套可复用的 BSP（板级支持包）”：从 GPIO 点亮 LED、读取按键，到串口、屏幕与 SPI Flash。掌握了这套 BSP 的组织方式，后续章节中的状态机、任务调度、菜单和协议代码就有了统一的“硬件底座”。

说明：后续例程会逐步统一到 HAL（Hardware Abstraction Layer）风格。HAL 本身可以由 STM32CubeMX 自动生成初始化代码，也可以由开发者手动编写，本附录展示的是“手动整理后的 HAL BSP”写法。

\medskip \textbf{本附录涉及到的实验内容：}
\begin{enumerate}
\item 开发板资源认识与 HAL GPIO 实验；
\item LED、KEY、BEEP 基础驱动；
\item HAL UART 重定向 printf；
\item TFT-LCD 显示驱动接口；
\item HAL SPI + W25Q64 存储驱动。
\end{enumerate}

% ---------------------------------------------------------------------
\section{单片机开发板——正点原子 STM32F1 战舰开发板}

本教程大部分硬件实验基于正点原子 STM32F1 战舰开发板。它采用 STM32F1，属于 Cortex-M3 内核：

\begin{itemize}
\item 主频最高 72MHz；
\item Flash 512KB、SRAM 64KB；
\item 丰富的 GPIO、定时器、USART、SPI、I2C、ADC 外设；
\item 板载 LED、按键、蜂鸣器、TFT-LCD 接口与 W25Q64 SPI Flash。
\end{itemize}

常用板载资源如下表。

\begin{table}[htbp]
\centering \caption{战舰开发板常用资源} \label{tab:board-resource}
    \begin{tabular}{lll}
        \toprule
        \textbf{资源} & \textbf{常用引脚} & \textbf{说明} \\
        \midrule
        LED0 与 LED1 & GPIO 输出脚 & HAL\_GPIO\_WritePin 控制 \\
        KEY0、KEY1 与 KEY2 & GPIO 输入脚 & HAL\_GPIO\_ReadPin 读取 \\
        BEEP & GPIO 输出脚 & 高低电平驱动蜂鸣器 \\
        USART1 & PA9 与 PA10 & HAL UART + printf \\
        TFT-LCD & FSMC 并行接口 & 显示菜单与状态 \\
        W25Q64 & SPI + 片选脚 & 存储参数、字库和数据 \\
        \bottomrule
    \end{tabular}
\end{table}

HAL 工程通常由 CubeMX 初始化时钟、引脚和中断，再在用户目录中补充应用逻辑。本附录展示的代码可以在 CubeMX 生成的工程基础上直接修改。

% ---------------------------------------------------------------------
\section{BSP 库的构建方式}

BSP 层负责把“芯片寄存器”封装成“面向应用的接口”。例如应用层不需要知道 LED 接在哪个引脚，只需要调用 \texttt{LED\_On()}。

参考作者实际工程的目录结构，建议把“统一外设初始化层”和“板级硬件驱动”分开：

\begin{lstlisting}
Board/System/
  bsp_periph.c   bsp_periph.h     /* GPIO/外设时钟/中断统一封装 */
Board/Hardware/
  bsp_led.c      bsp_led.h
  bsp_key.c      bsp_key.h
  bsp_beep.c     bsp_beep.h
  bsp_uart.c     bsp_uart.h
  bsp_lcd.c      bsp_lcd.h
  bsp_w25q64.c   bsp_w25q64.h
\end{lstlisting}

BSP 文件的编写原则：

\begin{enumerate}
\item 头文件只暴露必要的接口函数与宏；
\item \texttt{.c} 文件内部使用 \texttt{static} 隐藏硬件句柄细节；
\item 每个外设对应一个“初始化 + 操作”的模块；
\item 模块命名遵循“模块名 + 下划线 + 功能名”。
\end{enumerate}

在 HAL 中，每个外设通常有一个句柄，例如 \texttt{UART\_HandleTypeDef}、 \texttt{SPI\_HandleTypeDef}。BSP 内部持有这些句柄，并对外提供更贴近业务的 API。

为了让 GPIO 与时钟代码不被每个驱动重复书写，先建立统一的 \texttt{bsp\_periph} 层。它用 \texttt{GPIO\_IO\_T} 把“端口 + 引脚”合并成一个枚举，用 \texttt{GPIO\_MODE\_T} 表达引脚模式：

\begin{lstlisting}[style=c]
/* bsp_periph.h（节选）：PA0/PB12 等已由 GPIO_IO_T 统一编码 */
typedef enum __GPIO_MODE_T
{
    INPUT_U = 0,      /* 上拉输入 */
    INPUT_D = 1,      /* 下拉输入 */
    INPUT_F = 2,      /* 浮空输入 */
    OUTPUT_PP = 3,    /* 推挽输出 */
    OUTPUT_OD_U = 4,  /* 开漏输出 + 上拉 */
    OUTPUT_AF_PP = 7, /* 复用推挽输出 */
    ANALOG = 11       /* 模拟输入 */
} GPIO_MODE_T;

void BSP_GPIO_ClkEnable(GPIO_TypeDef *port);
void BSP_GPIO_Init(GPIO_IO_T io, GPIO_MODE_T mode);
void BSP_GPIO_InitEx(GPIO_IO_T io, GPIO_MODE_T mode,
                     uint8_t alternate);
void BSP_GPIO_WritePin(GPIO_IO_T io, uint8_t state);
uint8_t BSP_GPIO_ReadPin(GPIO_IO_T io);
void BSP_GPIO_TogglePin(GPIO_IO_T io);
void BSP_NVIC_Init(IRQn_Type irq, uint8_t preemption, uint8_t sub);
\end{lstlisting}

\texttt{BSP\_GPIO\_Init} 内部会自动调用 \texttt{BSP\_GPIO\_ClkEnable}，再把 \texttt{GPIO\_IO\_T} 与 \texttt{MODE\_T} 翻译成 HAL 的端口、引脚和模式。驱动里只需要写 \texttt{PA0}、\texttt{PB12}，不需要分别写 GPIOA + GPIO\_PIN\_0。

% ---------------------------------------------------------------------
\section{基础 GPIO 库的编写}

有了 \texttt{bsp\_periph} 后，驱动初始化不再直接写 HAL 结构体，而是调用统一入口：

\begin{lstlisting}[style=c]
#include "bsp_periph.h"

void Led_Init(void)
{
    BSP_GPIO_Init(PB12, OUTPUT_PP);   /* LED1 */
    BSP_GPIO_Init(PB0,  OUTPUT_PP);   /* LED2 */
}
\end{lstlisting}

其中 \texttt{bsp\_periph.c} 内部负责把模式枚举转换为 HAL 配置。常用模式：

\begin{itemize}
\item \texttt{OUTPUT\_PP}：推挽输出，驱动 LED、蜂鸣器等；
\item \texttt{INPUT\_U}：上拉输入，适合独立按键；
\item \texttt{INPUT\_D}：下拉输入；
\item \texttt{OUTPUT\_AF\_PP}：复用推挽输出，用于串口 TX 等。
\end{itemize}

STM32F1 没有独立复用功能号，传入 \texttt{GPIO\_AF\_NONE} 即可；F4 与 H7 上则传入 \texttt{GPIO\_AF7\_USART1} 这类复用功能号。

% ---------------------------------------------------------------------
\section{LED 代码框架构建}

LED 驱动文件先把引脚定义成宏，再调用 \texttt{bsp\_periph} 接口：

\begin{lstlisting}[style=c]
#define LED1_PIN  PB12
#define LED2_PIN  PB0

void LED_Init(void)
{
    BSP_GPIO_Init(LED1_PIN, OUTPUT_PP);
    BSP_GPIO_Init(LED2_PIN, OUTPUT_PP);
    BSP_GPIO_WritePin(LED1_PIN, GPIO_PIN_SET);  /* 高电平熄灭 */
    BSP_GPIO_WritePin(LED2_PIN, GPIO_PIN_SET);
}

/* 本开发板 LED 低电平点亮 */
void LED1_ON(void)
{
    BSP_GPIO_WritePin(LED1_PIN, GPIO_PIN_RESET);
}

void LED1_OFF(void)
{
    BSP_GPIO_WritePin(LED1_PIN, GPIO_PIN_SET);
}

void LED1_Turn(void)
{
    BSP_GPIO_TogglePin(LED1_PIN);
}
\end{lstlisting}

头文件只需要声明：

\begin{lstlisting}[style=c]
void LED_Init(void);
void LED1_ON(void);
void LED1_OFF(void);
void LED1_Turn(void);
void LED2_ON(void);
void LED2_OFF(void);
void LED2_Turn(void);
\end{lstlisting}

这样应用层不必关心引脚宏，后续换板时只改 BSP 源文件。

% ---------------------------------------------------------------------
\section{KEY 代码框架构建}

按键驱动不只提供“按下返回 1”，还应把物理电平抽象成按键电压枚举，方便上层状态机使用：

\begin{lstlisting}[style=c]
typedef enum __KeyVoltage_T
{
    KEY_VOLTAGE_RELEASED = 0,
    KEY_VOLTAGE_PRESSED
} KeyVoltage_T;

typedef enum __KeyEnum_T
{
    KEY_MODE,
    KEY_ENTER,
    KEY_UP,
    KEY_DOWN,
    KEY_NUMBER
} KeyEnum_T;

#define KEY_MODE_PIN  PA0
#define KEY_MODE_PULL INPUT_U

void Key_Init(void)
{
    BSP_GPIO_Init(KEY_MODE_PIN, KEY_MODE_PULL);
}

KeyVoltage_T Key_GetVoltage(KeyEnum_T key)
{
    switch (key) {
        case KEY_MODE:
            return BSP_GPIO_ReadPin(KEY_MODE_PIN) ?
                   KEY_VOLTAGE_PRESSED : KEY_VOLTAGE_RELEASED;
        default:
            return KEY_VOLTAGE_RELEASED;
    }
}
\end{lstlisting}

实际工程中，\texttt{bsp\_key.c} 会把每个按键的引脚\texttt{pin}、模式\texttt{mode}与触发电平\texttt{activeLevel} 放在配置数组里，并额外提供 \texttt{Key\_Tick()}、单击、双击和长按事件接口。简单轮询会产生抖动，因此本书把“按下、松开、单击和长按”等状态判断放进状态机章节处理；KEY 驱动只负责提供“当前电平” 这一底层事实，事件层交给状态机。

% ---------------------------------------------------------------------
\section{BEEP 代码框架构建}

蜂鸣器驱动与 LED 类似，使用同一个 GPIO 接口：

\begin{lstlisting}[style=c]
#define BEEP_GPIO_PIN PB1

void Beep_Init(void)
{
    BSP_GPIO_Init(BEEP_GPIO_PIN, OUTPUT_PP);
    BSP_GPIO_WritePin(BEEP_GPIO_PIN, GPIO_PIN_RESET); /* 默认关闭 */
}

void Beep_On(void)
{
    BSP_GPIO_WritePin(BEEP_GPIO_PIN, GPIO_PIN_SET);
}

void Beep_Off(void)
{
    BSP_GPIO_WritePin(BEEP_GPIO_PIN, GPIO_PIN_RESET);
}
\end{lstlisting}

实际工程还会提供 \texttt{Beep\_Play(onMs, offMs, cycles)} 与 \texttt{Beep\_Tick()}， \texttt{Beep\_Tick()} 通常放在定时器中断中按固定周期调用，由状态控制发声和静音节奏。如果希望发出不同音调，可以把 BEEP 引脚配置为定时器 PWM 输出，通过占空比和频率控制声音。

% ---------------------------------------------------------------------
\section{UART 代码框架构建}

HAL 串口使用 \texttt{UART\_HandleTypeDef} 句柄，但驱动文件不直接写 \texttt{\_\_HAL\_RCC\_USARTx\_CLK\_ENABLE}，而是调用 bsp\_periph 提供的 \texttt{BSP\_USART\_ClkEnable}：

\begin{lstlisting}[style=c]
static UART_HandleTypeDef s_Uart;   /* 句柄放在 .c 内部 */

void Uart_Init(void)
{
    /* 1. 使能时钟：bsp_periph 自动匹配 USART1 */
    BSP_USART_ClkEnable(USART1);

    /* 2. PA9=TX, PA10=RX：F1 复用模式无需 alternate */
    BSP_GPIO_InitEx(PA9, OUTPUT_AF_PP, GPIO_AF_NONE);
    BSP_GPIO_InitEx(PA10, OUTPUT_AF_PP, GPIO_AF_NONE);

    /* 3. 配置 USART1 */
    s_Uart.Instance = USART1;
    s_Uart.Init.BaudRate = 115200;
    s_Uart.Init.WordLength = UART_WORDLENGTH_8B;
    s_Uart.Init.StopBits = UART_STOPBITS_1;
    s_Uart.Init.Parity = UART_PARITY_NONE;
    s_Uart.Init.Mode = UART_MODE_TX_RX;
    HAL_UART_Init(&s_Uart);
}
\end{lstlisting}

发送一字节与 printf 重定向：

\begin{lstlisting}[style=c]
void Uart_SendByte(uint8_t byte)
{
    HAL_UART_Transmit(&s_Uart, &byte, 1u, HAL_MAX_DELAY);
}

int fputc(int ch, FILE *f)
{
    uint8_t byte = (uint8_t)ch;
    HAL_UART_Transmit(&s_Uart, &byte, 1u, HAL_MAX_DELAY);
    return ch;
}
\end{lstlisting}

重定向后：

\begin{lstlisting}[style=c]
printf("x=%d\r\n", x);
\end{lstlisting}

就能直接输出到串口助手。接收侧通常使用 \texttt{HAL\_UART\_Receive\_IT}，中断优先级统一通过 \texttt{BSP\_NVIC\_Init} 配置，完整设计见串口章节。

% ---------------------------------------------------------------------
\section{TFT-LCD 代码框架构建}

TFT-LCD 通常提供统一的“画点、写字符串和填充区域”接口。BSP 不需要让应用层直接操作控制器寄存器，而是提供：

\begin{lstlisting}[style=c]
void LCD_Init(void);
void LCD_Clear(uint16_t color);
void LCD_Fill(uint16_t x1, uint16_t y1,
              uint16_t x2, uint16_t y2,
              uint16_t color);
void LCD_ShowChar(uint16_t x, uint16_t y,
                  char ch, uint16_t size);
void LCD_ShowString(uint16_t x, uint16_t y,
                    char *str, uint16_t size);
\end{lstlisting}

菜单系统、状态机显示都依赖这组接口。只要 LCD 驱动符合上述函数签名，菜单层就与屏幕型号解耦。

% ---------------------------------------------------------------------
\section{W25Q64 代码框架构建}

W25Q64 是 8MB SPI Flash，用于保存参数、字库或用户数据。

\subsection{HAL SPI 代码构建}

HAL 的 SPI 同样使用句柄 \texttt{SPI\_HandleTypeDef}，时钟与 GPIO 通过 bsp\_periph 配置：

\begin{lstlisting}[style=c]
static SPI_HandleTypeDef s_Spi;

void Spi_Init(void)
{
    BSP_SPI_ClkEnable(SPI1);

    /* SPI1 常用引脚：PA5=SCK, PA6=MISO, PA7=MOSI */
    BSP_GPIO_InitEx(PA5, OUTPUT_AF_PP, GPIO_AF_NONE);
    BSP_GPIO_InitEx(PA6, OUTPUT_AF_PP, GPIO_AF_NONE);
    BSP_GPIO_InitEx(PA7, OUTPUT_AF_PP, GPIO_AF_NONE);

    s_Spi.Instance = SPI1;
    s_Spi.Init.Mode = SPI_MODE_MASTER;
    s_Spi.Init.Direction = SPI_DIRECTION_2LINES;
    s_Spi.Init.DataSize = SPI_DATASIZE_8BIT;
    s_Spi.Init.CLKPolarity = SPI_POLARITY_LOW;
    s_Spi.Init.CLKPhase = SPI_PHASE_1EDGE;
    s_Spi.Init.NSS = SPI_NSS_SOFT;
    s_Spi.Init.FirstBit = SPI_FIRSTBIT_MSB;
    s_Spi.Init.BaudRatePrescaler = SPI_BAUDRATEPRESCALER_16;
    HAL_SPI_Init(&s_Spi);
}

uint8_t Spi_SwapByte(uint8_t byte)
{
    uint8_t rxByte = 0;
    HAL_SPI_TransmitReceive(&s_Spi, &byte, &rxByte, 1u, HAL_MAX_DELAY);
    return rxByte;
}
\end{lstlisting}

片选引脚不参与 SPI 数据交换，由驱动单独用 GPIO 控制。在 STM32H743 的实际工程中，作者把 W25Q64 接到了 QSPI 外设，BSP 文件名为 \texttt{bsp\_qspi\_w25q64.c}，但对外接口依然是 “读、写和擦除”三段式。

\subsection{W25Q64 代码构建}

W25Q64 属于“命令 + 地址 + 数据”式 SPI 芯片。BSP 需要封装：

\begin{lstlisting}[style=c]
void W25Q64_Init(void);
uint8_t W25Q64_ReadID(void);           /* 读取厂商/型号 ID */
void W25Q64_EraseSector(uint32_t addr); /* 4KB 扇区擦除 */
void W25Q64_WritePage(uint32_t addr,
                      uint8_t *data,
                      uint16_t len);
void W25Q64_ReadData(uint32_t addr,
                     uint8_t *data,
                     uint16_t len);
\end{lstlisting}

读写之前必须通过片选引脚选中芯片，发送命令字节后再发送 24 位地址。写 Flash 前必须先擦除，因为 Flash 只能把 1 写成 0，不能把 0 写成 1。

% ---------------------------------------------------------------------
\section{小结}

BSP 的本质是“硬件隔离”：HAL 句柄和寄存器细节放在 .c 文件内部，头文件只暴露稳定 API。 LED、按键、蜂鸣器、串口、LCD 和 W25Q64 都遵守同一套框架后，后面的任务调度、菜单、协议和参数存储章节就不需要关心引脚细节，只需调用 BSP 接口。遇到新的 STM32 开发板时，也只需要替换 BSP 层，而应用层代码可以原封不动地保留。